\documentclass[11pt]{article}

\usepackage[a4paper,margin=1in]{geometry}
\usepackage{amsmath,amssymb,amsthm,mathtools}

\newtheorem{lemma}{Lemma}

\newcommand{\Db}{\mathrm{D}^{-}(R)}
\newcommand{\XX}{\mathbf{X}}
\newcommand{\TT}{\mathcal{T}}
\newcommand{\thicktensor}[1]{\operatorname{thick}^{\otimes}\{#1\}}

\title{Proof for Remark 7.22}
\author{ShuYu}
\date{}

\begin{document}

\maketitle

\section*{Problem Statement}

The main goal of this task is to investigate Remark 7.22; the rest is background. We do not know if it holds for $d \geqslant 3$, and this is the only question.

\medskip

\noindent\textbf{Example 7.21.}
Let $R$ be a commutative ring, and let $x \in R$ be a nonzerodivisor. For integers $a,b,c \geqslant 0$, define a complex
\[
X(a,b,c)=\bigoplus_{i \geqslant 0}(R/f_i)[i]
=\left(\cdots \xrightarrow{0} R/f_2 \xrightarrow{0} R/f_1 \xrightarrow{0} R/f_0 \to 0\right),
\]
where $f_i=x^{a i^2+b i+c} \in R$. Then
\[
\thicktensor{X(a,b,c)\mid a,b,c \geqslant 0}
=
\thicktensor{X(1,0,0)}.
\]

\medskip

\noindent\textbf{Remark 7.22.}
One can consider a general statement of Example 7.21 by defining
\[
f_i=x^{a_0 i^d+a_1 i^{d-1}+\cdots+a_d},
\]
so that Example 7.21 is exactly the case $d=2$. We do not know if it holds for $d \geqslant 3$.

\section*{Proof}

We prove that the general statement holds for all integers $d \geqslant 1$. Specifically, for any integer $d \geqslant 1$, any commutative noetherian ring $R$, and any nonzerodivisor $x \in R$, define
\[
X(a_0,a_1,\ldots,a_d)
=
\bigoplus_{i \geqslant 0}\left(R/x^{a_0 i^d + a_1 i^{d-1} + \cdots + a_d}\right)[i]
\]
for non-negative integers $a_0,\ldots,a_d$. We claim that
\[
\thicktensor{X(a_0,a_1,\ldots,a_d)\mid a_0,\ldots,a_d \geqslant 0}
=
\thicktensor{X(1,0,\ldots,0)}.
\]

The inclusion $(\supseteq)$ is trivial since $X(1,0,\ldots,0)$ appears on the left-hand side with $a_0=1$ and $a_1=\cdots=a_d=0$. For the inclusion $(\subseteq)$, we must show that every $X(a_0,\ldots,a_d)$ belongs to $\thicktensor{X(1,0,\ldots,0)}$. We abbreviate
\[
\TT=\thicktensor{X(1,0,\ldots,0)}
\]
throughout.

Let $e_k$ denote the $(d+1)$-tuple with $1$ in position $k$ and $0$ elsewhere, for $0 \leqslant k \leqslant d$. By Lemma 7.20 (second part, iterated), the complex $X(a_0,\ldots,a_d)$ whose exponent at index $i$ is $a_0 i^d + a_1 i^{d-1} + \cdots + a_d$ belongs to the thick closure of the complexes $X(e_k)$ for $k=0,\ldots,d$. Since $\operatorname{thick}^{\otimes}$ is at least as large as the thick closure, it suffices to show $X(e_k) \in \TT$ for each $k$.

Now $X(e_k)$ has exponent sequence $\{i^{d-k}\}_{i \geqslant 0}$. Define
\[
\XX(p)=\bigoplus_{i \geqslant 0}(R/x^{p(i)})[i]
\]
for any sequence $p(i) \geqslant 0$. Then our goal reduces to showing that for each $0 \leqslant k \leqslant d$, the complex $\XX(i^{d-k})$ belongs to $\TT=\thicktensor{\XX(i^d)}$. Equivalently, for each $0 \leqslant m \leqslant d$, we must show $\XX(i^m)\in\TT$.

The proof requires two new ingredients beyond those used in Example 7.21: a coefficient formula for finite differences, and a generalization of Proposition 7.2 / Corollary 7.3 from $2$-fold to $p$-fold splitting. We establish these first.

\subsection*{Lemma A (Coefficient Formula for Finite Differences)}

\begin{lemma}
For integers $d \geqslant 1$ and $0 \leqslant m \leqslant d$, the $m$-th forward finite difference of $i^d$ has the monomial expansion
\[
\Delta^m(i^d)
=
\sum_{j=0}^{d-m} \binom{d}{j}\, m!\, S(d-j,m)\, i^j,
\]
where $S(n,r)$ denotes the Stirling number of the second kind. In particular:
\begin{enumerate}
\item Every coefficient $\binom{d}{j}\, m!\, S(d-j,m)$ is a non-negative integer.
\item The leading coefficient, at $j=d-m$, equals $\dfrac{d!}{(d-m)!}$.
\item $\Delta^m(i^d)\geqslant 0$ for all $i \geqslant 0$.
\end{enumerate}
\end{lemma}

\begin{proof}
The $m$-th forward finite difference is
\[
\Delta^m(i^d)=\sum_{s=0}^{m}(-1)^{m-s}\binom{m}{s}(i+s)^d.
\]
Since the coefficient of $i^j$ in $(i+s)^d$ is $\binom{d}{j}s^{d-j}$, the coefficient of $i^j$ in $\Delta^m(i^d)$ is
\[
\binom{d}{j}\sum_{s=0}^{m}(-1)^{m-s}\binom{m}{s}s^{d-j}.
\]
By the standard identity for Stirling numbers of the second kind,
\[
\sum_{s=0}^{m}(-1)^{m-s}\binom{m}{s}s^n=m!\,S(n,m)
\]
for all $n \geqslant 0$, with the convention $S(0,0)=1$ and $S(n,m)=0$ for $n<m$. Therefore the coefficient of $i^j$ is exactly $\binom{d}{j}\,m!\,S(d-j,m)$.

For (1), each factor $\binom{d}{j}$, $m!$, and $S(d-j,m)$ is a non-negative integer. For (2), at $j=d-m$, the coefficient is
\[
\binom{d}{d-m}m!S(m,m)
=
\frac{d!}{(d-m)!\,m!}\cdot m!\cdot 1
=
\frac{d!}{(d-m)!}.
\]
For $j>d-m$, we have $d-j<m$, so $S(d-j,m)=0$; thus the polynomial has degree exactly $d-m$. For (3), since
\[
\Delta^m(i^d)=\sum_j \alpha_j i^j
\]
with all $\alpha_j \geqslant 0$, and $i \geqslant 0$, we get $\Delta^m(i^d)\geqslant 0$.
\end{proof}

\subsection*{Lemma B ($p$-fold Splitting, Generalized Corollary 7.3)}

\begin{lemma}
Let $p \geqslant 2$ be an integer, and let
\[
X=\bigoplus_{n \geqslant 0} M_n[n]
\]
be a complex in $\Db$ with zero differentials. For $r=0,1,\ldots,p-1$, define
\[
X^{(r)}=\bigoplus_{i \geqslant 0} M_{pi+r}[i].
\]
Then
\[
X \in \thicktensor{X^{(0)},X^{(1)},\ldots,X^{(p-1)}}.
\]
\end{lemma}

\begin{proof}
We show that for each $r \in \{0,\ldots,p-1\}$, the complex
\[
Z_r:=\bigoplus_{i \geqslant 0} M_{pi+r}[pi+r]
\]
is in $\thicktensor{X^{(r)}}$. Since
\[
X=Z_0 \oplus Z_1 \oplus \cdots \oplus Z_{p-1}
\]
as a direct sum decomposition of graded modules with zero differentials, the conclusion follows.

Fix $r \in \{0,\ldots,p-1\}$. Define the complex
\[
Y_r=\bigoplus_{j \geqslant 0} R[(p-1)j+r].
\]
Under the convention that $[k]$ places the module in homological degree $k$ (and thus cohomological degree $-k$), $Y_r$ has non-zero components only in cohomological degrees $\leqslant 0$, so $Y_r \in \Db$.

Because $X^{(r)}$ is a complex with zero differentials and $Y_r$ is a complex of free $R$-modules, hence $K$-flat, the derived tensor product $X^{(r)} \otimes_R^{L} Y_r$ equals the ordinary tensor product and has zero differentials. Explicitly,
\[
X^{(r)} \otimes_R^{L} Y_r
=
\bigoplus_{i,j \geqslant 0}(M_{pi+r}\otimes_R R)[i+(p-1)j+r]
=
\bigoplus_{i,j \geqslant 0} M_{pi+r}[i+(p-1)j+r].
\]

In homological degree $n$, this tensor product has module
\[
\bigoplus_{i+(p-1)j+r=n} M_{pi+r}.
\]
Consider the diagonal subcollection defined by $j=i$. For $j=i$, the degree is
\[
i+(p-1)i+r=pi+r,
\]
and the module contributed is $M_{pi+r}$. Thus, for each $i_0 \geqslant 0$, the module $M_{pi_0+r}$ appears as a direct summand of the degree-$(pi_0+r)$ component of $X^{(r)} \otimes_R^{L} Y_r$.

Since all differentials are zero, the subcomplex
\[
Z_r=\bigoplus_{i \geqslant 0} M_{pi+r}[pi+r]
\]
is a direct summand of $X^{(r)} \otimes_R^{L} Y_r$ in the category of complexes, hence in $\Db$. Since thick tensor ideals are closed under tensor product with any object of $\Db$, we have
\[
X^{(r)} \otimes_R^{L} Y_r \in \thicktensor{X^{(r)}}.
\]
Because thick tensor ideals are closed under direct summands, we conclude $Z_r \in \thicktensor{X^{(r)}}$.

Therefore
\[
X=\bigoplus_{r=0}^{p-1} Z_r \in \thicktensor{X^{(0)},X^{(1)},\ldots,X^{(p-1)}}.
\]
\end{proof}

\subsection*{Lemma C (Subtraction of Exponent Sequences)}

\begin{lemma}
If the complexes $\XX(a_i+b_i)$ and $\XX(b_i)$ both belong to $\TT$, and $a_i,b_i \geqslant 0$ for all $i \geqslant 0$, then $\XX(a_i)$ also belongs to $\TT$.
\end{lemma}

\begin{proof}
For each $i \geqslant 0$, we have a short exact sequence of $R$-modules
\[
0 \to R/x^{a_i} \xrightarrow{x^{b_i}} R/x^{a_i+b_i} \to R/x^{b_i} \to 0,
\]
where the map is well-defined and injective since $x$ is a nonzerodivisor. Taking the direct sum over all $i \geqslant 0$ and placing each term in degree $[i]$, we obtain a short exact sequence of complexes with zero differentials:
\[
0 \to \XX(a_i) \to \XX(a_i+b_i) \to \XX(b_i) \to 0.
\]
Since $\TT$ is a thick subcategory, closed under extensions and cones, and both $\XX(a_i+b_i)$ and $\XX(b_i)$ are in $\TT$, the two-out-of-three property gives $\XX(a_i)\in\TT$.
\end{proof}

\subsection*{Lemma D (Finite Differences Preserve the Thick Tensor Ideal)}

\begin{lemma}
Let $\{a_i\}_{i \geqslant 0}$ be a sequence of non-negative integers such that $a_{i+1} \geqslant a_i$ for all $i \geqslant 0$. If $\XX(a_i)\in\TT$, then $\XX(a_{i+1}-a_i)\in\TT$.
\end{lemma}

\begin{proof}
For each $i \geqslant 0$, the short exact sequence
\[
0 \to R/x^{a_i} \xrightarrow{x^{a_{i+1}-a_i}} R/x^{a_{i+1}} \to R/x^{a_{i+1}-a_i} \to 0
\]
yields, upon taking direct sums over $i$, a short exact sequence of complexes with zero differentials:
\[
0 \to \XX(a_i) \to \XX(a_{i+1}) \to \XX(a_{i+1}-a_i) \to 0.
\]
By the two-out-of-three property for thick subcategories, it suffices to show that $\XX(a_{i+1})\in\TT$.

By re-indexing $j=i+1$, we have
\[
\XX(a_{i+1})
=
\bigoplus_{i \geqslant 0}(R/x^{a_{i+1}})[i]
=
\bigoplus_{j \geqslant 1}(R/x^{a_j})[j-1]
=
\left(\bigoplus_{j \geqslant 1}(R/x^{a_j})[j]\right)[-1].
\]
Since
\[
\XX(a_i)=\bigoplus_{j \geqslant 0}(R/x^{a_j})[j]
\]
is a complex with zero differentials, it decomposes as a direct sum of its graded pieces, so
\[
\bigoplus_{j \geqslant 1}(R/x^{a_j})[j]
\]
is a direct summand of $\XX(a_i)\in\TT$, and hence belongs to $\TT$. Applying the shift $[-1]$ gives $\XX(a_{i+1})\in\TT$.
\end{proof}

\paragraph{Corollary of Lemma D.}
By iterating Lemma D, if $\XX(a_i)\in\TT$ and the sequence $\{a_i\}$ satisfies $\Delta^l(a_i)\geqslant 0$ for all $i \geqslant 0$ and $l=1,\ldots,m$, then $\XX(\Delta^m(a_i))\in\TT$. In particular, since $\Delta^l(i^d)\geqslant 0$ for all $i \geqslant 0$ and $0 \leqslant l \leqslant d$ by Lemma A, starting from $\XX(i^d)\in\TT$ we get $\XX(\Delta^m(i^d))\in\TT$ for every $0 \leqslant m \leqslant d$.

\subsection*{Proof of the Main Claim}

We prove by induction on $k$, from $k=0$ to $k=d$, that $\XX(i^k)\in\TT$.

\paragraph{Base case ($k=0$).}
The complex $\XX(1)$ has exponent sequence $\{1,1,1,\ldots\}$. Since
\[
\XX(i^d)=\bigoplus_{i \geqslant 0}(R/x^{i^d})[i]
\]
has $(R/x^{1^d})[1]=(R/x)[1]$ as a direct summand, the module $R/x$ belongs to $\TT$. Then
\[
\XX(1)=(R/x)\otimes_R^{L}\left(\bigoplus_{i \geqslant 0} R[i]\right)\in\TT
\]
since $\bigoplus_{i \geqslant 0} R[i]\in\Db$ and thick tensor ideals are closed under tensor products.

\paragraph{Inductive step.}
Suppose $k \geqslant 1$ and $\XX(i^l)\in\TT$ for all $0 \leqslant l<k$. We prove $\XX(i^k)\in\TT$. Set $m=d-k$. By the corollary of Lemma D, the complex $\XX(Q(i))\in\TT$, where
\[
Q(i)=\Delta^m(i^d).
\]
By Lemma A,
\[
Q(i)=\sum_{j=0}^{k}\alpha_j i^j,
\qquad
\alpha_j=\binom{d}{j}m!S(d-j,m)\geqslant 0,
\]
and
\[
\alpha_k=\frac{d!}{k!}.
\]

Define
\[
c=\alpha_k=\frac{d!}{k!},
\qquad
q(i)=Q(i)-c\,i^k=\sum_{j=0}^{k-1}\alpha_j i^j.
\]

\paragraph{Step 1.}
We show $\XX(q(i))\in\TT$. Since each $\alpha_j$ is a non-negative integer by Lemma A, and $\XX(i^j)\in\TT$ for $j<k$ by the inductive hypothesis, Lemma 7.20 (scalar multiplication) gives $\XX(\alpha_j i^j)\in\TT$. Summing with Lemma 7.20 (addition), we get $\XX(q(i))\in\TT$.

\paragraph{Step 2.}
We show $\XX(c\,i^k)\in\TT$. Since $Q(i)=c\,i^k+q(i)$ with $c\,i^k\geqslant 0$ and $q(i)\geqslant 0$, Lemma C applied with $a_i=c\,i^k$ and $b_i=q(i)$ gives $\XX(c\,i^k)\in\TT$.

\paragraph{Step 3.}
We show $\XX(i^k)\in\TT$ via $c$-fold splitting. Apply Lemma B with $p=c$ to the complex $\XX(i^k)$. The $c$-fold parts are
\[
\XX(i^k)^{(r)}=\XX((ci+r)^k),
\qquad
r=0,\ldots,c-1.
\]
By Lemma B,
\[
\XX(i^k)\in\thicktensor{\XX((ci+r)^k)\mid r=0,\ldots,c-1}.
\]
It remains to show each $\XX((ci+r)^k)\in\TT$.

By the binomial theorem,
\[
(ci+r)^k=c^k i^k+\sum_{j=0}^{k-1}\binom{k}{j}c^j r^{k-j} i^j.
\]
For the lower-order terms, each coefficient
\[
\beta_j=\binom{k}{j}c^j r^{k-j}
\]
is a non-negative integer. By the inductive hypothesis, $\XX(i^j)\in\TT$, so $\XX(\beta_j i^j)\in\TT$, and their sum is in $\TT$. For the leading term, Step 2 gives $\XX(c\,i^k)\in\TT$. By Lemma 7.20 (scalar multiplication with factor $c^{k-1}$),
\[
\XX(c^k i^k)=\XX(c^{k-1}\cdot c\,i^k)\in\TT.
\]
Adding the leading and lower-order terms using Lemma 7.20 (addition), we get $\XX((ci+r)^k)\in\TT$. By Lemma B, this implies $\XX(i^k)\in\TT$.

\subsection*{Completing the Proof}

By the main claim, $\XX(i^m)=X(e_{d-m})\in\TT$ for all $0 \leqslant m \leqslant d$. By Lemma 7.20 (second part, iterated), for any tuple $(a_0,\ldots,a_d)$ of non-negative integers, the complex $X(a_0,\ldots,a_d)$ with exponent
\[
\sum_j a_j i^{d-j}
\]
belongs to the thick closure of $\{X(e_0),X(e_1),\ldots,X(e_d)\}$, which is contained in $\TT$. Therefore
\[
\thicktensor{X(a_0,\ldots,a_d)\mid a_0,\ldots,a_d \geqslant 0}
\subseteq
\TT
=
\thicktensor{X(1,0,\ldots,0)},
\]
proving the general statement.

\section*{Key Ideas}

\begin{enumerate}
\item \textbf{$p$-fold splitting (Lemma B).} This generalizes Corollary 7.3 to arbitrary $p \geqslant 2$, enabling division by $p$ for any $p$. It is proved by tensoring with $Y_r=\bigoplus_j R[(p-1)j+r]$ and extracting the diagonal direct summand.
\item \textbf{Coefficient formula via Stirling numbers (Lemma A).} The finite difference $\Delta^{d-k}(i^d)$ is a degree-$k$ polynomial whose lower-degree terms have non-negative integer coefficients, enabling subtraction through Lemma C.
\item \textbf{Bottom-up induction.} Starting from $\XX(i^d)\in\TT$, we iterate finite differences to produce lower-degree polynomials, subtract off strictly lower-degree terms already known to lie in $\TT$, and use $c$-fold splitting to isolate $\XX(i^k)$.
\end{enumerate}

\end{document}
